\documentclass[8pt]{article}


\usepackage[T2A]{fontenc}
\usepackage[utf8]{inputenc}
\usepackage[english,russian]{babel}
\usepackage{blindtext}
\usepackage[a4paper, total={8.0in, 10in}]{geometry}

\usepackage[utf8]{inputenc}


\usepackage{amsmath,amssymb,amsfonts,amsthm}
\usepackage{mathrsfs}  
\usepackage{cancel}
\usepackage{leftindex}  

\usepackage{xcolor}
\setcounter{MaxMatrixCols}{20}
% \usepackage{lscape}

\newcommand{\bo}[1] {\mathbf{#1}}
\newcommand{\R} {\mathbb{R}}
\newcommand{\norm}[1] {\left|\left|{#1}\right|\right|}
\newcommand{\T}{^\top}



% \newcommand{\cos}[1] {\text{cos}\left(#1\right)}
% \newcommand{\sin}[1] {\text{cos}\left(#1\right)}



\newtheorem{assumption}{Assumption}
\newtheorem{remark}{Remark}
\newtheorem{theorem}{Theorem}


\definecolor{myhotpink}{RGB}{255, 80, 200}
\definecolor{mywarmpink}{RGB}{255, 60, 160}
\definecolor{mylightpink}{RGB}{255, 80, 200}
\definecolor{mypalepink}{RGB}{255, 200, 230}
\definecolor{mypink}{RGB}{255, 30, 80}
\definecolor{mydarkpink}{RGB}{155, 25, 60}

\definecolor{mypaleblue}{RGB}{240, 240, 255}
\definecolor{mylightblue}{RGB}{120, 150, 255}
\definecolor{myblue}{RGB}{90, 90, 255}
\definecolor{mygblue}{RGB}{70, 110, 240}
\definecolor{mydarkblue}{RGB}{0, 0, 180}
\definecolor{myblackblue}{RGB}{40, 40, 120}

\definecolor{mygreen}{RGB}{0, 200, 0}
\definecolor{mygreen2}{RGB}{245, 255, 230}

\definecolor{mygray}{gray}{0.8}
\definecolor{mydarkgray}{RGB}{80, 80, 160}
\definecolor{mylightgray}{RGB}{160, 160, 160}
\definecolor{mylightgreyred}{RGB}{210, 160, 160}



\definecolor{mydarkred}{RGB}{160, 30, 30}
\definecolor{mylightred}{RGB}{255, 150, 150}
\definecolor{myred}{RGB}{200, 110, 110}
\definecolor{myblackred}{RGB}{120, 40, 40}

\definecolor{myred1}{RGB}{100, 0, 0}
\definecolor{myred2}{RGB}{160, 60, 60}
\definecolor{myred3}{RGB}{220, 90, 120}
\definecolor{myred4}{RGB}{250, 120, 120}
\definecolor{myred5}{RGB}{255, 150, 150}

\definecolor{mygreen}{RGB}{0, 200, 0}
\definecolor{mygreen2}{RGB}{205, 255, 200}
\definecolor{mydarkgreen}{RGB}{0, 100, 0}

\definecolor{mydarkcolor}{RGB}{60, 25, 155}
\definecolor{mylightcolor}{RGB}{130, 180, 250}



\newcommand{\re}[1] {\textcolor{mydarkpink}{#1}}

%%%%%%%%%%%%%%%%%%%%%%%%%%%%%%%%%%%%%%%

\usepackage{tcolorbox}
\tcbuselibrary{skins,breakable}
\usetikzlibrary{shadings,shadows}

\newenvironment{myexampleblock}[1]{%
    \tcolorbox[beamer,%
    noparskip,breakable,
    colback=mylightgreygreen,colframe=myblackgreen,%
    colbacklower=mylime!75!mylightgreygreen,%
    title=#1]}%
    {\endtcolorbox}

\newenvironment{myalertblock}[1]{%
    \tcolorbox[beamer,%
    noparskip,breakable,
    colback=mylightgreyred,colframe=myblackred,%
    colbacklower=myred!75!mylightgreyred,%
    title=#1]}%
    {\endtcolorbox}

\newenvironment{myblock}[1]{%
    \tcolorbox[beamer,%
    noparskip,breakable,
    colback=mylightblue,colframe=myblackblue,%
    colbacklower=myblackblue!75!myblue,%
    title=#1]}%
    {\endtcolorbox}

\title{MyMathTemplate}


\begin{document}

\section*{Домашняя работа 3 }
Рассмотрим линеаризованную модель квадрокоптера, работающего в режиме висения (горизонтальный полёт). Представление в пространстве состояний имеет вид:
\[
\dot{\mathbf{x}} = A\mathbf{x} + B\mathbf{u}
\]
где вектор состояния \(\mathbf{x} \in \mathbb{R}^{12}\) и управляющий вход \(\mathbf{u} \in \mathbb{R}^4\) определяются следующим образом:

\[
\begin{aligned}
x_1 &= x \quad &\text{(инерциальная x-координата)} \\
x_2 &= y \quad &\text{(инерциальная y-координата)} \\
x_3 &= z \quad &\text{(инерциальная z-координата)} \\
x_4 &= \phi \quad &\text{(угол крена)} \\
x_5 &= \theta \quad &\text{(угол тангажа)} \\
x_6 &= \psi \quad &\text{(угол рысканья)} \\
x_7 &= \dot{x} \quad &\text{(скорость по x)} \\
x_8 &= \dot{y} \quad &\text{(скорость по y)} \\
x_9 &= \dot{z} \quad &\text{(скорость по z)} \\
x_{10} &= \dot{\phi} = p \quad &\text{(угловая скорость крена)} \\
x_{11} &= \dot{\theta} = q \quad &\text{(угловая скорость тангажа)} \\
x_{12} &= \dot{\psi} = r \quad &\text{(угловая скорость рысканья)}
\end{aligned}
\]

\[
\mathbf{u} = \begin{bmatrix} u_1 \\ u_2 \\ u_3 \\ u_4 \end{bmatrix} = \begin{bmatrix} \Delta\tau_\phi \\ \Delta\tau_\theta \\ \Delta\tau_\psi \\ \Delta T \end{bmatrix} \quad \text{(момент крена, момент тангажа, момент рысканья, изменение тяги)}
\]

Матрицы системы имеют вид:

\[
A = 
\begin{bmatrix}
0 & 0 & 0 & 0 & 0 & 0 & 1 & 0 & 0 & 0 & 0 & 0 \\
0 & 0 & 0 & 0 & 0 & 0 & 0 & 1 & 0 & 0 & 0 & 0 \\
0 & 0 & 0 & 0 & 0 & 0 & 0 & 0 & 1 & 0 & 0 & 0 \\
0 & 0 & 0 & 0 & 0 & 0 & 0 & 0 & 0 & 1 & 0 & 0 \\
0 & 0 & 0 & 0 & 0 & 0 & 0 & 0 & 0 & 0 & 1 & 0 \\
0 & 0 & 0 & 0 & 0 & 0 & 0 & 0 & 0 & 0 & 0 & 1 \\
0 & 0 & 0 & 0 & g & 0 & 0 & 0 & 0 & 0 & 0 & 0 \\
0 & 0 & 0 & -g & 0 & 0 & 0 & 0 & 0 & 0 & 0 & 0 \\
0 & 0 & 0 & 0 & 0 & 0 & 0 & 0 & 0 & 0 & 0 & 0 \\
0 & 0 & 0 & 0 & 0 & 0 & 0 & 0 & 0 & 0 & 0 & 0 \\
0 & 0 & 0 & 0 & 0 & 0 & 0 & 0 & 0 & 0 & 0 & 0 \\
0 & 0 & 0 & 0 & 0 & 0 & 0 & 0 & 0 & 0 & 0 & 0
\end{bmatrix}
\qquad
B = 
\begin{bmatrix}
0 & 0 & 0 & 0 \\
0 & 0 & 0 & 0 \\
0 & 0 & 0 & 0 \\
0 & 0 & 0 & 0 \\
0 & 0 & 0 & 0 \\
0 & 0 & 0 & 0 \\
0 & 0 & 0 & 0 \\
0 & 0 & 0 & 0 \\
0 & 0 & 0 & 1/m \\
1/I_{xx} & 0 & 0 & 0 \\
0 & 1/I_{yy} & 0 & 0 \\
0 & 0 & 1/I_{zz} & 0
\end{bmatrix}
\]


с параметрами \(g = 9.81\,\text{м/с}^2\), \(m = 1.0\,\text{кг}\), \(I_{xx} = I_{yy} = 0.01\,\text{кг·м}^2\), \(I_{zz} = 0.02\,\text{кг·м}^2\).

\section*{Задание}

Найдите закон управления по обратной связи по состоянию \(\mathbf{u} = K\mathbf{x}\), который асимптотически стабилизирует квадрокоптер. Используйте метод Ляпунова и сформулируйте задачу как полуопределённую программу (SDP).

    
\subsection*{Результаты:}
    \begin{itemize}
        \item Запишите ЛМН и код для соответствующей SDP.
        \item Предоставьте результирующую матрицу коэффициентов усиления \(K\).
        \item Промоделируйте реакцию замкнутой системы из начального условия \(\mathbf{x}(0) = [1, 0, 1, 0.1, 0.1, 0, 0, 0, 0, 0, 0, 0]^T\) и постройте графики 12 переменных состояния за 10 секунд.
        \item (Дополнительно) Сгенерируйте 100 различных и случайных начальных состояний и рассчитайте траектории движения квадракоптера. Найдите долю устойчивых траекторий.
\end{itemize}


\end{document}